\section{Related work}\label{sec:related}

\paragraph{Soft Thinking and noise.}
\citet{zhang2025soft} introduced \ST{}, a training-free method that feeds probability-weighted mixtures of token
embeddings during thinking, and reported small accuracy gains and shorter generations on math and code benchmarks.
\citet{wu2025single} probed the method and found that models rely mainly on the top token of the blend, which
produces the Greedy Pitfall; they call the models single-threaded reasoners.
They proposed Stochastic \ST{} with Dirichlet or Gumbel-Softmax noise \citep{jang2017categorical,maddison2017concrete}
and motivated the Gumbel variant through the Gumbel-max trick and its top-$k$ extension \citep{kool2019stochastic}.
They did not isolate the weight off the top token in a controlled test.
We make the consequence of the trick explicit, check it in logged runs, and isolate the residual with a coupled
ablation, which refines their single-threaded account for the Gumbel variant.
Noisy soft tokens are also used for training: \citet{butt2025soft} add noise to soft input embeddings for
reinforcement-learning exploration, and \citet{zheng2025softgrpo} inject Gumbel noise and use Gumbel-Softmax in
policy optimisation.
Our analysis concerns a fixed, untrained model at inference.

\paragraph{Is superposition used?}
\citet{zhu2025superposition} prove that continuous thoughts can encode several search frontiers in one state, and
\citet{hao2024coconut} report breadth-first-search-like latent states in Coconut.
\citet{deng2025latent} train latent reasoning to act as a superposition of reasoning paths.
\citet{rizvimartel2026illusion} test whether superposition occurs in three regimes and find signs of it only in
models trained from scratch with latent thoughts; in the training-free regime, vanilla \ST{} looks much like
discrete decoding under the logit lens.
We study the training-free regime with noise, condition on forks where the model is uncertain, and go beyond
the next-token readout to storage in the key-value cache and to causal use.

\paragraph{Collapse of mixtures.}
\citet{backour2026dynamics} injects two-way mixtures of embeddings into pretrained models and projects each layer's
state onto the line between the two pure states.
Backour finds that trained models collapse mixtures far more than random-initialised controls and that softmax readout
with autoregressive feedback amplifies small imbalances until one component dominates.
Our snap score is his CALIB measure applied at natural forks of a reasoning model.
We then ask what the collapse leaves behind, what blends the noise feeds, and whether the residual changes
predictions and answers.

\paragraph{Causal analysis of latent reasoning.}
Activation patching localises information by swapping internal states between runs \citep{meng2022locating}; our
key-value patching applies it to the cache entries of one position.
For trained latent reasoners, \citet{li2026dynamics} intervene on latent steps of Coconut and CODI and report a gap
between early output bias and late representational commitment.
\citet{zhang2025latent} find that Coconut's latent tokens are insensitive to steering and support shortcuts, and
\citet{aswal2026observable} show that decodable latent patterns also appear in controls and need not affect behaviour.
\citet{jin2026final} track latent-reasoning faithfulness across training.
Following their advice, each of our causal tests has a matched control: a control token in Q2 and a random residual
in Q4.

\paragraph{Two senses of superposition.}
In interpretability, superposition usually means that a network represents more features than it has dimensions
\citep{elhage2022toy}.
We use the word only in the latent-reasoning sense above: several candidate continuations carried and used at once.
